\documentclass[../main]{subfiles}
\begin{document}
\ifSubfilesClassLoaded{\setcounter{page}{4}}{}
\section{Comunismo e coscienza sociale}
\label{sec:10_comunismo_coscienza}

Secondo la prospettiva marxista, è attraverso la produzione materiale che la società genera la propria coscienza. L'umanità ha conosciuto diverse forme di coscienza sociale, sviluppatesi nel corso della storia. A partire dalla coscienza mitologica, in cui l'ideale non è ancora nettamente distinto, o in cui, a seconda delle fasi, non lo è del tutto, dall'attività materiale immediata, con l'avvento della divisione del lavoro e dell'alienazione dell'uomo dal proprio lavoro, emergono la filosofia e la religione, e successivamente l'arte e la scienza. Contestualmente, si sviluppano la morale e il diritto, necessari quando la società richiede una regolamentazione esterna dei rapporti tra individui con interessi contrastanti. La coscienza sociale non è semplicemente la somma delle coscienze individuali, bensì la coscienza della società nel suo complesso, distinta dall'attività umana immediata, la sua dimensione ideale.Allo stesso modo, la coscienza individuale non è altro che la coscienza sociale, interiorizzata dall’individuo in relazione al suo ruolo e alla sua posizione all’interno della società.

La coscienza sociale, nelle sue molteplici forme, assume un’esistenza oggettiva nei confronti dei membri della comunità. Per poter vivere in una società umana, con le sue esigenze e i suoi principi morali, giuridici, religiosi, scientifici e persino artistici e filosofici, gli individui sono tenuti a conformarsi ad essa, così come devono adattarsi all’ambiente circostante e al livello di progresso tecnologico. Nella coscienza sociale si riflette e si esprime non un elemento arbitrario, bensì il modo in cui l’uomo interagisce con la natura e con i suoi simili. La coscienza sociale, in tutte le sue manifestazioni, non è un prodotto diretto della natura, ma il risultato della sua appropriazione da parte della società umana, una modalità di produzione che si riflette in essa.

L'idea che la coscienza sociale si emancipi dalle attività pratiche, assumendo forme di esistenza autonome, e che imponga ai membri della società un vincolo morale, è stata riconosciuta ben prima di Marx. Essa costituisce il fondamento non solo della visione religiosa del mondo, ma anche di tutta la filosofia idealista.

L'avvento del pensiero comunista ha richiesto una critica sia di questo fenomeno in sé, sia delle modalità con cui esso viene interpretato nel pensiero teorico. La critica della coscienza sociale, non solo nel contesto di un'epoca specifica, ma in riferimento alla società in cui prevale lo sfruttamento dell'uomo sull'uomo, ha svolto un ruolo cruciale nella definizione del contenuto ideologico del comunismo.

Come ben si sa, il marxismo affonda le sue radici nella critica alla filosofia hegeliana e, più specificamente, nell’analisi di Feuerbach della religione cristiana, forma di coscienza sociale che ha dominato l’Europa nel Medioevo e che, pur persistendo, continua a rivendicare un ruolo preminente anche nella società borghese. L’importanza che Marx attribuiva alla critica di Feuerbach alla religione è evidente nel modo in cui egli la considera, definendola un’impresa a tutti gli effetti. Ma in cosa consiste, secondo Marx, questa “impresa”, in relazione alla domanda che ci siamo posti?«1) nella dimostrazione che la filosofia non è altro che la religione espressa in forma di pensiero e sistematicamente organizzata, una mera trasfigurazione, un diverso modo di manifestarsi dell’alienazione dell’essenza umana, e che, perciò, anch’essa è destinata all’obsolescenza (questa è un’osservazione di notevole importanza, concernente la filosofia in quanto forma di coscienza sociale); 2) nell’affermazione del vero materialismo e della scienza autentica, giacché Feuerbach fa del rapporto sociale “uomo-uomo” il principio cardine della sua teoria.»

Questi sono i punti di forza del pensiero di Feuerbach, ma la sua coerenza non è assoluta, come Marx ebbe modo di evidenziare: «Feuerbach parte dall’analisi dell’alienazione religiosa, dalla scissione del mondo in una sfera religiosa, idealizzata, e una sfera reale. Si dedica quindi a ricondurre la sfera religiosa alle sue fondamenta terrene. Tuttavia, non si rende conto che, una volta compiuto questo passo, resta ancora molto da fare. In particolare, il fatto che queste fondamenta terrene si separino da sé stesse e si proiettino in un regno autonomo può essere spiegato unicamente attraverso l’autodistruzione e l’autocontraddittorietà insite in tali fondamenta. Ne consegue che queste ultime devono essere prima comprese nella loro intrinseca contraddizione e, successivamente, trasformate in modo concreto mediante l’eliminazione di tale contraddizione».Ne consegue che, una volta risolto il mistero della sacra famiglia, ad esempio nell’ambito della famiglia terrena, quest’ultima deve essere sottoposta a un’analisi critica e trasformata in modo radicale nella prassi. La vita terrena dell’uomo, la sua attività produttiva, diviene così, per Marx, l’oggetto di studio volto a comprendere l’essenza e le radici terrene della coscienza. «La religione è essa stessa priva di contenuto; le sue origini non risiedono nel cielo, ma sulla terra, e con la distruzione di quella realtà distorta, di cui essa rappresenta l’espressione teorica, essa si estingue di per sé» (Marx a Arnold Ruge, Dresda, 30 novembre 1842).

Lo stesso principio si applica a tutte le altre forme di coscienza sociale: la filosofia, la morale, il diritto, e persino la scienza e l’arte – già citate – devono essere considerate come manifestazioni della prassi umana, derivanti dalla sua essenza e dal suo contenuto intrinseco, e non come entità separate e indipendenti. All’interno del pensiero marxista, si traccia una chiara distinzione tra la base delle relazioni sociali – intesa come l’insieme delle forze produttive e dei rapporti di produzione – e la sovrastruttura che da essa scaturisce, e alla quale appartiene, appunto, la coscienza sociale.

In uno dei suoi appunti preparatori per *Il Capitale*, Marx affronta la questione del diritto, sostenendo: «Se queste banalità – ovvero la concezione del diritto propria degli economisti contemporanei – fossero ridotte alla loro essenza, rivelerebbero più di quanto i loro stessi fautori possano immaginare. In particolare, esse implicano che ogni forma di produzione dà origine alle relazioni giuridiche e alle strutture di governo ad essa intrinsecamente connesse. L’errore e la mancanza di perspicacia risiedono nel fatto che fenomeni organicamente legati vengono considerati in una relazione puramente casuale, come se fossero collegati da un nesso meramente razionale. Gli economisti borghesi, infatti, si limitano a considerare che, nell’ambito del sistema di governo vigente, sia possibile ottenere una maggiore efficienza produttiva rispetto a quanto si potrebbe raggiungere, ad esempio, attraverso il ricorso alla forza bruta. Essi dimenticano, tuttavia, che anche la forza bruta rappresenta una forma di diritto e che il diritto del più forte persiste, seppur in una veste diversa, persino nel loro presunto “stato di diritto”.»

È necessario non solo comprendere, ma anche trasformare le «radici terrestri» della coscienza sociale: occorre mutarne la natura, riformare il carattere della produzione sociale. Il fatto che le forme della coscienza sociale derivino dal modo di produzione sociale dimostra che la loro essenza è ben più complessa di quanto possa apparire a prima vista. Per esempio, l’eliminazione completa del diritto, il superamento della sua necessità come regolatore esterno, è possibile solo a determinate condizioni sociali, in una società comunista sviluppata e fondata su solide basi. Lo stesso principio vale per la religione: non si può semplicemente abolire, dichiarandola fallace. È necessario superarla, preservando al contempo tutti i risultati della cultura umana che si sono sviluppati in forma religiosa, e distruggendo, attraverso una trasformazione radicale, il carattere delle attività umane svolte quotidianamente.

Certamente, in un breve scritto dedicato al tema della coscienza sociale nel contesto del comunismo, possiamo affrontare la questione solo a grandi linee. Tuttavia, è fondamentale sottolineare un aspetto cruciale riguardante la società che precede il comunismo, una società intrinsecamente divisa in classi. Analogamente a come gli individui occupano posizioni distinte all'interno del processo produttivo, in relazione alla loro appartenenza di classe, anche la coscienza sociale assume una specifica connotazione di classe. Ciò non implica l'inesistenza di un ambito di coscienza unitario, né la negazione di una verità, di un bene e di una bellezza oggettivi, che trascendano la mera posizione di classe del loro sostenitore. Piuttosto, si tratta di riconoscere che tale ambito di coscienza rappresenta esso stesso un campo di battaglia tra le classi. In altre parole, la classe economicamente dominante estende il proprio predominio anche alla sfera della coscienza, e la coscienza di tale classe si afferma come egemone all'interno della società stratificata.Sulla scia del pensiero di Marx e Engels, Lenin evidenziò l’importanza cruciale della battaglia sul piano ideologico, ovvero dello scontro teorico, concepito come una peculiare forma di lotta di classe.

Nella brochure intitolata «Cosa fare?», Lenin dedicò un intero paragrafo all’importanza della battaglia teorica condotta dal proletariato, soffermandosi anche sulle riflessioni di Engels in merito. Oggi, ci appare particolarmente perspicace la seguente osservazione leniniana: «Chiunque abbia una conoscenza, seppur minima, della realtà del nostro movimento non può fare a meno di notare come l’ampia diffusione del marxismo sia stata accompagnata da un certo appiattimento del livello teorico. Molti individui, scarsamente o per nulla preparati sul piano teorico, si sono avvicinati al movimento attratti dal suo valore pratico e dai suoi successi concreti. Di conseguenza, si può facilmente comprendere l’incongruenza che denota l’atteggiamento del «Rab. Delo» quando, con tono trionfalistico, riporta l’affermazione di Marx: «Ogni passo compiuto dal movimento reale vale più di una dozzina di programmi».In un'epoca segnata dalla confusione teorica, ripetere queste parole equivale a gridare «non potete più cambiare dimora!» durante un corteo funebre. Queste parole di Marx provengono infatti dalla sua lettera sul Programma di Gotha, nella quale egli critica con veemenza l'eclettismo che aveva caratterizzato la formulazione dei principi: se si desiderava giungere a una sintesi, scriveva Marx ai leader del partito, allora era necessario stringere dei patti per perseguire gli obiettivi pratici del movimento, ma era inammissibile mercanteggiare sui principi, non si doveva scendere a compromessi sul piano teorico. Questa era l'intenzione di Marx, eppure ci sono individui che, invocando il suo nome, si adoperano per sminuire l'importanza della teoria! Senza una teoria rivoluzionaria, non può esistere un movimento rivoluzionario. Non si può sottolineare abbastanza questo concetto, in un momento storico in cui la moderna predicazione dell'opportunismo si accompagna a un'adesione alle forme più restrittive di azione pratica.

Ricollegandosi all'eredità di Engels, Lenin riporta una lunga citazione che riteniamo opportuno trascrivere anche qui: «I lavoratori tedeschi vantano due pregi fondamentali rispetto ai loro omologhi del resto d'Europa. In primo luogo, appartengono alla nazione più incline alla riflessione teorica del continente e hanno preservato in sé quella predisposizione che sembra essersi quasi del tutto estinta nelle cosiddette classi “colte” della Germania. Senza il contributo della filosofia tedesca, in particolare quella di Hegel, il socialismo scientifico tedesco – l’unico esempio di socialismo scientifico degno di questo nome – non si sarebbe mai sviluppato. E senza la sensibilità teorica dei lavoratori, questo socialismo scientifico non avrebbe mai potuto penetrare così profondamente nelle loro coscienze, come possiamo constatare oggi.»L’entità di questo vantaggio è evidente: da un lato, l’indifferenza verso ogni teoria rappresenta uno dei principali ostacoli al progresso del movimento operaio inglese, nonostante l’eccellente organizzazione delle singole corporazioni; dall’altro, il proudonismo, nella sua forma originaria, ha generato confusione e instabilità tra i francesi e i belgi, mentre, nella sua versione caricaturale, ridotta a tal quale da Bakunin, ha prodotto gli stessi effetti tra gli spagnoli e gli italiani.

Un secondo vantaggio risiede nel fatto che i tedeschi si sono uniti al movimento operaio quasi in ultimo luogo. Così come il socialismo teorico tedesco non deve mai dimenticare di fondarsi sulle idee di Saint-Simon, Fourier e Owen – tre pensatori che, al di là di ogni idealismo e utopismo, si distinguono per la loro straordinaria acutezza intellettuale e la cui capacità di anticipare innumerevoli verità, oggi dimostriamo con rigore scientifico – allo stesso modo, il movimento operaio tedesco, nella sua prassi, non deve mai dimenticare di essersi sviluppato sulle spalle dei movimenti inglese e francese, potendo ora beneficiare della loro preziosa esperienza e, conseguentemente, evitare gli errori che, in quel contesto, erano per lo più inevitabili.In che stato saremmo oggi senza l’esempio offerto dai sindacati inglesi e dalla battaglia politica intrapresa dai lavoratori francesi, senza quell’immenso slancio che, in particolare, emanò dalla Comune di Parigi?

È doveroso riconoscere che i lavoratori tedeschi hanno saputo sfruttare con straordinaria perizia i vantaggi derivanti dalla loro posizione. Per la prima volta, nella storia del movimento operaio, la lotta si è sviluppata in modo sistematico in tutti i suoi tre ambiti, in una dinamica coordinata e interconnessa: quello teorico, quello politico e quello pratico-economico (la resistenza contro i capitalisti).

Questo brano mette in luce l'atteggiamento di Engels nei confronti della battaglia ideologica. Tuttavia, la lotta di classe nel dominio della coscienza sociale non si esaurisce nella sola sfera teorica, così come la coscienza sociale non si limita alla pura elaborazione concettuale. La contesa si estende anche al campo dell'attività estetica, a quello della morale e si confronta con il diritto e la religione. Anche tale aspetto non va sottovalutato. Nondimeno, l'impegno teorico volto a favorire l'acquisizione e lo sviluppo, da parte del proletariato, della propria coscienza di classe rappresenta la chiave di volta per tutte le altre forme di consapevolezza, in quanto consente di comprenderle appieno.

È doveroso sottolineare come, all’interno di una società di classe, l’interesse universale si manifesti, a livello della coscienza collettiva, attraverso la classe stessa, indipendentemente dalla sua effettiva comprensione da parte degli individui. La coscienza della classe rivoluzionaria, che in questa particolare fase storica incarna l’interesse universale, costituisce la vera forza motrice della storia, e proprio attraverso di essa tale interesse viene realizzato. In quest’ottica, non solo la morale, la religione e la filosofia, ma anche il diritto, la scienza e l’arte si rivelano forme di ideologia, poiché sia l’estetica nell’arte, intesa come modalità di appropriazione del mondo, sia la conoscenza scientifica sono condizionate – e, di conseguenza, anche limitate – dal movimento storico della società e dalla configurazione delle forze in gioco nella lotta di classe.

A tal proposito, mi preme affrontare il tema del soggetto del pensiero in Marx e Lenin. È doveroso ricordare come Marx abbia distinto l'economia politica in due filoni, quello classico e quello volgare, sulla base di un preciso criterio. Egli evidenziò che l'economia politica classica rappresentò una fase cruciale nell'evoluzione della comprensione del proprio oggetto di studio, in contrasto con la sua controparte volgare. Definì inequivocabilmente la natura parziale, o meglio, la sua "classicità", di tale disciplina. Sia l'economia politica classica che quella volgare, peraltro, restano scienze borghesi, funzionali agli interessi della borghesia ed espressione, a livello teorico, delle sue direttrici di sviluppo. Nel momento in cui tale sviluppo smette di coincidere con quello della società nel suo complesso, cessa di esserne una fedele rappresentazione; in altri termini, quando la borghesia smette di essere una classe rivoluzionaria, l'economia politica si trasforma inevitabilmente in un discorso "volgare".Ciò implica che non si tratta più di un’idea in evoluzione, bensì di un vero e proprio freno all’avanzamento del pensiero.

Chiaramente, la questione non concerne le qualità individuali dei singoli studiosi di scienze politiche. Marx dimostra, in modo inequivocabile, che il soggetto del pensiero non è rappresentato né dai singoli teorici, né dalla società nella sua interezza, ma da una classe sociale. E, inoltre, tale classe è il soggetto del pensiero solo in quanto è anche il soggetto dell’azione storica.

Ritroviamo un’eco di questa stessa idea nelle opere di Lenin, laddove egli distingue tra due categorie di partiti nel contesto della filosofia: il partito che favorisce il progresso del pensiero e quello che, al contrario, gli si oppone. Ciò deriva dal fatto che un partito incarna una classe sociale, la quale rappresenta il motore dell’azione storica e, di conseguenza, anche il soggetto del pensiero; l’altro partito, invece, non possiede questa caratteristica. Questa dinamica non si limita al solo ambito filosofico, ma si estende a tutte le discipline scientifiche, comprese quelle naturali. Sebbene possa sembrare meno evidente, ad esempio, ottenere finanziamenti per progetti di ricerca destinati a produrre risultati solo nel lunghissimo periodo è diventato un’impresa ardua, al punto che si può interpretare questa difficoltà come una manifestazione, seppur mascherata, della lotta di classe. Questa interpretazione resta valida fintanto che lo sviluppo della società è determinato dalla dinamica della lotta di classe. Possiamo quindi affermare che la verità oggettiva è intrinsecamente orientata e presenta una chiara connotazione di classe. Un discorso analogo può essere applicato all’oggettività dei concetti di bene e di bellezza.

Nel settembre del 1917 (agosto secondo il calendario giuliano), Lenin formulò un pensiero che, a nostro avviso, riveste un'importanza capitale e che si pone in perfetta continuità con quanto precedentemente esposto. Tale pensiero illustra il modo in cui, in un'epoca rivoluzionaria, il proletariato sviluppa e affina la propria coscienza di classe. È fondamentale considerare che l'obiettivo del proletariato, a differenza di quello perseguito dalle altre classi, non consiste affatto nell'ottenere e consolidare la propria egemonia, bensì nell'abolire, una volta raggiunta, ogni forma di dominio di classe e ogni divisione della società in classi. Lenin affermò con fermezza: «Saremo intransigenti nel sottoporre a vaglio anche i più piccoli dubbi, avvalendoci del giudizio dei lavoratori più consapevoli, del giudizio del nostro partito, in cui riponevamo la nostra fiducia, e del quale vedevamo l'intelligenza, l'onore e la coscienza del nostro tempo; nell'unione internazionale dei rivoluzionari internazionalisti scorgiamo l'unica garanzia per il movimento di liberazione della classe lavoratrice».«L’egemonia di una classe è inconcepibile senza la sua predominanza nell’ambito della coscienza», affermava Lenin. Per trionfare, la classe rivoluzionaria deve confrontarsi, nel dominio della coscienza, non con parametri esterni, bensì con i propri, definiti dalla sua specifica missione storica. Una volta sottratta alla borghesia la pretesa di egemonia in tale ambito, e avvalendosi delle più avanzate conquiste derivanti dalla sua coscienza di classe, il proletariato si trasforma in una forza materiale e concreta di trasformazione storica: «L’arma della critica non può certo sostituire la critica delle armi; la forza materiale deve essere sopraffatta dalla forza materiale. Tuttavia, la teoria si rivela una forza materiale nel momento stesso in cui si appropria delle masse», scriveva Marx nella «Critica della filosofia hegeliana del diritto».

Fino ad ora, tuttavia, abbiamo affrontato il tema della coscienza sociale non in relazione al comunismo, bensì nel contesto della situazione attuale, che costituisce il punto di partenza per la transizione verso di esso. Qual è, dunque, l'accento posto dal comunismo, in quanto scienza, sulla coscienza sociale?

Il comunismo pone l'accento sul superamento della divisione del lavoro all'interno della società, sull'eliminazione dell'alienazione dell'uomo rispetto agli altri uomini, al processo e ai risultati della sua attività. Questo implica anche i corrispondenti mutamenti nell'ambito della coscienza sociale. L'alienazione delle diverse forme della coscienza sociale, sia l'una rispetto all'altra, sia direttamente rispetto all'attività umana, sarà anch'essa superata.

Le capacità ideali e operative dell'uomo, attualmente «limitate» all'interno delle forme della coscienza sociale, e considerate ambiti separati dalla produzione materiale – come l'attività estetica, morale e intellettuale – diventeranno di pubblico dominio.

La problematica della coscienza sociale nel contesto comunista si configura come una questione eminentemente pratica. Essa concerne la pianificazione e la sua implementazione nell’ambito dell’attività materiale, in modo tale da indurre i mutamenti corrispondenti anche nella sfera spirituale. Attraverso la trasformazione delle proprie azioni, gli individui modificano, di conseguenza, la propria coscienza, intesa quale espressione ideale di tale attività. Tuttavia, nel comunismo, la stessa trasformazione della prassi assume un carattere consapevole. Di conseguenza, non è più la mera esistenza sociale a determinare la coscienza sociale, bensì quest’ultima, concepita come la consapevolezza di un’associazione libera e solidale di individui, plasma l’esistenza, che a sua volta influenza e definisce tale coscienza.

La collettivizzazione dell’intero patrimonio culturale e spirituale, e la sua riorganizzazione secondo una logica coerente, diviene un’esigenza imprescindibile per la transizione verso il comunismo. Pertanto, all’interno della teoria comunista, le questioni relative alla coscienza sociale acquistano significato solo se interpretate come aspetti intrinsecamente legati alla costruzione del comunismo.

Nell’ambito dell’attività comunista, a differenza del lavoro alienato, l’ideale – inteso come ideale in senso stretto e in relazione al “mondo delle idee” – finisce per coincidere, in ultima analisi e a livello empirico. L’ideale, dunque, si manifesta sia come riflesso (rappresentazione, espressione) del materiale nel materiale stesso, ovvero come riflesso attivo della realtà oggettiva, sia come ideale obiettivo dell’azione, a cui quest’ultima deve tendere, sia come ideale estetico. In ciò risiede l’elemento cruciale per superare la divisione del lavoro e la sua alienazione in tutte le sue forme, affermando così l’autovalore intrinseco dell’attività, indipendentemente dai risultati concreti che essa produce.

È fondamentale precisare che sarebbe un errore considerare gli aspetti consapevoli e materiali dell'attività umana come aventi lo stesso valore nell'ambito del lavoro, a prescindere dalle sue specifiche caratteristiche storiche, e a maggior ragione quando si considera l'attività della società nel suo complesso. Naturalmente, l'aspetto ideale è presente in ogni attività come elemento costitutivo della definizione degli obiettivi, a meno che tale definizione non sia esterna all'attività stessa (nel caso di un lavoratore che offre semplicemente la propria forza lavoro, i suoi obiettivi non coincidono affatto con quelli della produzione), ma qui si tratta di una questione diversa. In ogni caso, il modo in cui viene prodotta la vita materiale determina il modo in cui viene prodotta la dimensione ideale. E questa è una realtà ineludibile con cui è necessario confrontarsi.

Questo assume una rilevanza pratica non trascurabile, poiché, sebbene un’idea acquisisca una forza tangibile nel momento in cui permea le masse, essa stessa è il frutto di una profonda comprensione delle dinamiche materiali che ne regolano l’azione (ne rappresenta, al contempo, il risultato e una delle manifestazioni dell’interazione tra il mondo materiale). In altre parole, per mutare il processo di generazione delle idee, è necessario intervenire sul processo di produzione della vita materiale. Solo in questo modo, e non nel senso opposto, si potrà ottenere il cambiamento desiderato. Un approccio inverso risulterebbe inefficace.

Per quanto potremmo impegnarci a diffondere (o, più precisamente, a implementare in maniera compiuta) la dialettica, essa non si trasformerebbe per questo in un modo di pensare diffuso, e la produzione del pensiero dialettico rimarrebbe, per così dire, un’attività elitaria. Questa attività, indubbiamente, riveste una grande importanza e da essa dipenderà in larga misura il futuro. Tuttavia, è necessario comprendere che coloro che adottano un approccio dialettico al pensiero non costituiranno una massa indistinta, date le attuali modalità di produzione della vita materiale e le corrispondenti modalità di formazione della coscienza.

Per quanto concerne la questione del comunismo, non basta limitarsi a osservare la direzione in cui l'esistenza materiale dell'uomo si evolve verso la coscienza – il dilemma del "quale viene prima" – ma è imperativo proporre soluzioni concrete per realizzare tale trasformazione. In tale contesto, le singole idee in merito assumono un'importanza cruciale. Sono fondamentali le proposte per la creazione di specifiche forme empiriche di comunismo, poiché non tutti gli aspetti materiali si trasformano istantaneamente in coscienza, ma lo fanno sempre attraverso modalità ben definite. Hanno, ad esempio, un valore significativo concetti quali quelli sviluppati da Makarenko: l'amore è un sentimento universale, che deve essere adeguatamente organizzato.

Qualunque importanza si possa attribuire al momento più propizio per l’azione, esso resta sempre legato alla vita e all’attività dell’individuo concreto, immerso nel mondo materiale. L’ideale rappresenta, in ultima analisi, una caratteristica intrinseca all’attività pratica dell’individuo sociale, e ne riflette lo sviluppo. Nel “regno della libertà”, ovvero nel comunismo compiuto, si distingue, tra le altre peculiarità, la capacità di elevare la coscienza sociale a un livello collettivo, attraverso un processo di produzione condivisa.

Già oggi si intravedono alcuni progressi nel campo della coscienza sociale, progressi che trascendono i confini del capitalismo. Ne abbiamo discusso, tra l’altro, nel secondo capitolo del presente lavoro. Si tratta dell’emergere di un nuovo ideale, plasmato nel contesto delle tecnologie moderne. Ma non si limita a questo: non ogni ideale è intrinsecamente legato alla coscienza sociale. Non è più una mera attività idealistica, confinata nella sua alienazione e caratterizzata da una natura sovra-strutturale, né un ideale relegato entro i ristretti limiti di un cretinismo professionale. Si tratta piuttosto del momento ideale che presiede alla produzione della vita umana in generale. E il modo in cui questa vita viene prodotta – un processo che, come nel caso dei database e del software libero, presenta delle analogie – determina la genesi di una nuova forma di coscienza sociale, immediata e non alienata, sebbene al momento circoscritta a una scala locale.Non si tratta semplicemente di creare un programma, ma piuttosto di raggiungere una forma di consapevolezza che scaturisce dalla stessa essenza del processo creativo e dal suo effettivo funzionamento.

La coscienza sociale trova la sua definizione solo nel concetto stesso di esistenza sociale. Quest’ultima non va intesa come entità separata dall’individuo nella sua interezza, né come pura astrazione ideale, bensì come una fase dell’attività dell’individuo concreto nel mondo materiale. Si tratta, in definitiva, della manifestazione consapevole dell’individuo nella sua totalità: un soggetto che agisce, percepisce e riflette in una dimensione universale. Tale attività, indipendentemente dalla sua natura specifica, si concretizzerà nella realizzazione di una forma diretta e collettiva del comunismo, divenendo un processo orientato alla creazione della libertà, alla definizione di obiettivi condivisi, all’affinamento della sensibilità, alla promozione della moralità e allo sviluppo del pensiero.

In sostanza, ciò significa che, nell’ambito del comunismo, ogni azione deve necessariamente comprendere anche la produzione di idee e la formulazione di un ideale per la futura esistenza, contribuendo alla creazione dell’uomo stesso.
\end{document}
